\chapter{Kafka and RabbitMQ Internals}
\label{ch:kafka-rabbitmq}

\begin{chaptergoalsbox}
\textbf{What you will learn:}
\begin{itemize}[leftmargin=12pt]
  \item Kernel I/O fundamentals from first principles: User Space vs. Kernel Space, OS Page Cache, and traditional \texttt{read()}/\texttt{write()} (4 copies) vs. \texttt{sendfile(2)} zero-copy.
  \item Apache Kafka storage and replication: segmented commit logs, sparse \texttt{.index} binary search, ISR, High Watermark (HW), and KRaft (\texttt{\_\_cluster\_metadata}).
  \item Kafka producer batching (\texttt{linger.ms}, \texttt{batch.size}, Kafka 3.0+ idempotent defaults) and consumer group rebalancing (\texttt{Range}, \texttt{RoundRobin}, \texttt{CooperativeSticky}).
  \item RabbitMQ AMQP 0-9-1 routing (Direct, Fanout, Topic, Headers exchanges), Erlang BEAM reductions, Mnesia metadata, and Quorum Queues vs. Classic Mirrored Queues.
\end{itemize}

\textbf{What you should already understand:}
\begin{itemize}[leftmargin=12pt]
  \item Raft majority quorums from Chapter~\ref{ch:cap-consistency} and queues vs. streams from Chapter~\ref{ch:messaging-systems}.
  \item No prior Linux kernel or Erlang background is assumed---we introduce page caches, system calls, and BEAM scheduling directly.
\end{itemize}

\textbf{Why this topic exists:}
Apache Kafka and RabbitMQ look similar on architecture slides---both move messages between services---but their internals are opposites. Treating them interchangeably leads to disaster: storing long-lived audit streams inside RabbitMQ RAM, or trying to route millions of per-user priority tasks through Kafka partitions.

\textbf{Motivating practical problem:}
A logistics platform uses RabbitMQ Classic Mirrored Queues to dispatch 40,000 shipment updates per second. When one broker reboots for security patching, resynchronizing mirrored queues locks the broker for four minutes and stalls checkout. After migrating to Kafka without understanding consumer rebalancing, deploying a new pod triggers a 90-second ``stop-the-world'' rebalance storm every time autoscaling runs.
\end{chaptergoalsbox}

\section{Kernel I/O Primer: Page Cache and Zero-Copy}

Before examining Kafka's disk format, we must answer a puzzle: how does Kafka sustain gigabytes per second of throughput when it writes every message to disk rather than keeping everything in application RAM? The secret lives inside the Linux operating system kernel.

\subsection{User Space, Kernel Space, and the OS Page Cache}

Modern CPUs enforce hardware privilege rings separating memory into two regions:
\begin{itemize}[leftmargin=18pt]
  \item \textbf{Kernel Space}: Privileged OS territory that owns hardware drivers (disks, network cards), TCP sockets, and the Virtual File System (VFS).
  \item \textbf{User Space}: Unprivileged memory where your Python, Java, or Go application runs. An application cannot touch a disk or network card directly; it must ask the kernel via a \textbf{system call} (such as \texttt{read()}, \texttt{write()}, or \texttt{sendfile()}), which pauses the thread and switches CPU mode into the kernel (a \textbf{context switch}).
\end{itemize}

Whenever the kernel reads or writes a file on disk, it automatically caches recently used $4\text{ KB}$ disk blocks inside free physical RAM called the \textbf{OS Page Cache}. Crucially, Kafka runs on the JVM, where caching millions of message objects inside Java heap memory wastes RAM on object headers and triggers long Stop-The-World garbage collection pauses. Instead, Kafka writes raw sequential bytes to files and lets the Linux Page Cache hold hot data in RAM with zero JVM GC overhead.

\subsection{Traditional \texttt{read()}/\texttt{write()} vs. \texttt{sendfile(2)} Zero-Copy}

Hardware controllers use \textbf{DMA (Direct Memory Access)} to transfer bytes between Disk/NIC and Kernel RAM without burning CPU cycles. Watch what happens when a traditional server reads a file from disk and sends it over a socket using standard \texttt{read()} followed by \texttt{write()} (Figure~\ref{fig:v4-ch05-zero-copy-comparison}, top):
\begin{enumerate}[leftmargin=18pt]
  \item \texttt{read()} context-switches to kernel; DMA copies bytes from Disk $\to$ OS Page Cache (Copy~1).
  \item CPU copies bytes from Page Cache $\to$ User-Space buffer (Copy~2); context-switches back to user space.
  \item \texttt{write()} context-switches to kernel; CPU copies bytes from User-Space buffer $\to$ Kernel Socket Buffer (Copy~3).
  \item Context-switches back to user space while DMA copies Socket Buffer $\to$ NIC (Copy~4).
\end{enumerate}
That costs \textbf{4 context switches and 4 data copies} (two of them wasting CPU) even though the broker never modified a single byte!

\begin{figure}[htbp]
\centering
\begin{tikzpicture}[
  scale=0.82,
  box/.style={rectangle, draw=headingblue, fill=codebg, thick, rounded corners=2pt, minimum width=2.3cm, minimum height=0.95cm, align=center, font=\sffamily\tiny}
]
  % Top: Traditional 4-copy
  \node[anchor=west, font=\sffamily\bfseries\footnotesize, color=warningRed] at (0, 2.3) {(a) Traditional \texttt{read()} + \texttt{write()}: 4 Context Switches, 4 Copies (2 CPU Copies)};
  \node[box] (d1) at (1.0, 1.3) {Disk};
  \node[box, fill=exampleBg] (pc1) at (4.2, 1.3) {OS Page Cache\\(Kernel)};
  \node[box, fill=warningBg] (usr1) at (7.4, 1.3) {App Buffer\\(User Space)};
  \node[box, fill=exampleBg] (sk1) at (10.6, 1.3) {Socket Buffer\\(Kernel)};
  \node[box] (nic1) at (13.8, 1.3) {NIC};

  \draw[-{Stealth[scale=0.9]}, thick] (d1) -- node[above, font=\sffamily\tiny] {1. DMA} (pc1);
  \draw[-{Stealth[scale=0.9]}, thick, color=warningRed] (pc1) -- node[above, font=\sffamily\tiny] {2. CPU Copy} (usr1);
  \draw[-{Stealth[scale=0.9]}, thick, color=warningRed] (usr1) -- node[above, font=\sffamily\tiny] {3. CPU Copy} (sk1);
  \draw[-{Stealth[scale=0.9]}, thick] (sk1) -- node[above, font=\sffamily\tiny] {4. DMA} (nic1);

  % Bottom: Zero-copy sendfile
  \node[anchor=west, font=\sffamily\bfseries\footnotesize, color=intuitionTeal] at (0, -0.3) {(b) Linux \texttt{sendfile(2)} Zero-Copy: 2 Context Switches, 2 DMA Copies (0 CPU Copies)};
  \node[box] (d2) at (1.0, -1.3) {Disk};
  \node[box, fill=intuitionBg] (pc2) at (5.0, -1.3) {OS Page Cache\\(Kernel)};
  \node[box, fill=intuitionBg] (sk2) at (9.2, -1.3) {Socket Buffer\\(Descriptor Only)};
  \node[box] (nic2) at (13.8, -1.3) {NIC};

  \draw[-{Stealth[scale=0.9]}, thick] (d2) -- node[above, font=\sffamily\tiny] {1. DMA} (pc2);
  \draw[-{Stealth[scale=0.9]}, thick, dashed, color=linkblue] (pc2) -- node[above, font=\sffamily\tiny] {Pointer + Length} (sk2);
  \draw[-{Stealth[scale=0.9]}, thick, color=intuitionTeal] (pc2) to [bend right=24] node[below, font=\sffamily\tiny] {2. DMA Scatter-Gather directly to NIC} (nic2);
  \draw[-{Stealth[scale=0.9]}, thick] (sk2) -- (nic2);
\end{tikzpicture}
\caption{Side-by-Side Comparison: Traditional 4-Copy Path vs. Kafka \texttt{sendfile(2)} Zero-Copy}
\label{fig:v4-ch05-zero-copy-comparison}
\end{figure}

Kafka replaces \texttt{read()}+\texttt{write()} with the Linux \texttt{sendfile(2)} syscall (Figure~\ref{fig:v4-ch05-zero-copy-comparison}, bottom). Only a tiny descriptor (memory offset and byte count) is placed in the Socket Buffer, and the NIC reads directly from the OS Page Cache via DMA scatter-gather. Data never crosses into JVM user space---cutting CPU copies to \textbf{zero} (note: TLS encryption in user space prevents kernel zero-copy unless kernel TLS, \texttt{kTLS}, is enabled).

\section{Apache Kafka Internals}

Designed at LinkedIn in 2010, Kafka models messaging as a distributed, partitioned, replicated append-only commit log.

\subsection{Commit Log Segments and Sparse Index}

A Kafka \textbf{Topic} is split into ordered \textbf{Partitions}. On disk, each partition directory stores rolling $1\text{ GB}$ \textbf{segment files}:
\begin{lstlisting}[language=bash]
/var/lib/kafka/data/orders-0/
    00000000000000000000.log        # Records starting at offset 0
    00000000000000000000.index      # Offset -> byte position sparse map
    00000000000000000000.timeindex  # Timestamp -> offset sparse map
\end{lstlisting}

\begin{figure}[htbp]
\centering
\begin{tikzpicture}[
  scale=0.85,
  entry/.style={rectangle, draw=exampleBlue, fill=exampleBg, thick, minimum width=1.1cm, minimum height=0.7cm, font=\sffamily\scriptsize}
]
  \node[anchor=west, font=\sffamily\bfseries\footnotesize, color=headingblue] at (0, 2.2) {Closed Segment 0000000000.log};
  \node[entry] (e0) at (0.6, 1.4) {Off: 0};
  \node[entry] (e1) at (1.8, 1.4) {Off: 1};
  \node[entry] (e2) at (3.0, 1.4) {Off: 2};
  \node[entry] (e3) at (4.2, 1.4) {Off: 3};

  \node[anchor=west, font=\sffamily\bfseries\footnotesize, color=headingblue] at (6.5, 2.2) {Active Segment 0000000004.log};
  \node[entry, fill=intuitionBg] (e4) at (7.1, 1.4) {Off: 4};
  \node[entry, fill=intuitionBg] (e5) at (8.3, 1.4) {Off: 5};
  \node[entry, fill=gray!15] (e6) at (9.5, 1.4) {Append};

  \draw[thick, dashed, color=gray!60] (0, 0.4) -- (11, 0.4);
  \node[anchor=west, font=\sffamily\bfseries\footnotesize, color=headingblue] at (0, 0.0) {Sparse Memory-Mapped Index (.index)};

  \node[rectangle, draw=hoodSlate, fill=hoodBg, font=\sffamily\tiny, align=center] at (1.8, -0.8) {Offset 0 $\to$ Pos 0};
  \node[rectangle, draw=hoodSlate, fill=hoodBg, font=\sffamily\tiny, align=center] at (4.8, -0.8) {Offset 4096 $\to$ Pos 65536};
  \node[rectangle, draw=hoodSlate, fill=hoodBg, font=\sffamily\tiny, align=center] at (8.2, -0.8) {Offset 8192 $\to$ Pos 131072};
\end{tikzpicture}
\caption{Kafka Segmented Commit Log and Sparse Memory-Mapped Index}
\label{fig:v4-ch05-kafka-segments}
\end{figure}

Rather than indexing every single message, Kafka writes an entry to \texttt{.index} roughly every $4\text{ KB}$ (\texttt{index.interval.bytes}). To fetch offset $4{,}120$, Kafka binary-searches the memory-mapped \texttt{.index} to find the nearest entry $\le 4{,}120$ (offset $4{,}096$ at byte $65{,}536$) and scans forward sequentially a few kilobytes inside the \texttt{.log} file.

\subsection{Replication Topology: ISR, High Watermark, and KRaft}

Every partition has one \textbf{Leader} broker handling reads/writes and zero or more \textbf{Follower} brokers pulling updates:
\begin{itemize}[leftmargin=18pt]
  \item \textbf{In-Sync Replicas (ISR)}: Followers caught up within \texttt{replica.lag.time.max.ms} (default $30\text{ s}$).
  \item \textbf{Log End Offset (LEO)}: Offset of the next record to be appended on a replica.
  \item \textbf{High Watermark (HW)}: The minimum LEO across all brokers in the ISR. Consumers are only allowed to read up to the High Watermark so they never observe uncommitted records that could vanish on leader failover.
\end{itemize}

\begin{figure}[htbp]
\centering
\begin{tikzpicture}[
  scale=0.85,
  nodebox/.style={rectangle, draw=headingblue, fill=codebg, thick, rounded corners=2pt, minimum width=3.2cm, minimum height=1.2cm, align=center, font=\sffamily\footnotesize}
]
  \node[nodebox, fill=exampleBg] (leader) at (0, 0) {\textbf{Leader (Broker 1)}\\LEO = 10\\HW = 8};
  \node[nodebox, fill=intuitionBg] (f1) at (5, 1.2) {\textbf{Follower 1 (Broker 2)}\\LEO = 9\\In ISR};
  \node[nodebox, fill=intuitionBg] (f2) at (5, -1.2) {\textbf{Follower 2 (Broker 3)}\\LEO = 8\\In ISR};

  \draw[-{Stealth[scale=1.0]}, thick] (f1) -- node[above, font=\sffamily\tiny] {Fetch RPC} (leader);
  \draw[-{Stealth[scale=1.0]}, thick] (f2) -- node[below, font=\sffamily\tiny] {Fetch RPC} (leader);

  \node[anchor=west, font=\sffamily\footnotesize, color=headingblue] at (8.5, 0) {High Watermark = $\min(\text{LEO}_{\text{ISR}}) = 8$};
\end{tikzpicture}
\caption{Kafka ISR Replication and High Watermark Calculation}
\label{fig:v4-ch05-kafka-hw}
\end{figure}

Historically, Kafka stored cluster topology and partition leaders in Apache ZooKeeper, capping clusters around $\sim 200{,}000$ partitions due to ZooKeeper memory reload times on failover. Under \textbf{KRaft} (Kafka Raft Metadata Mode, KIP-500), ZooKeeper is removed entirely: dedicated controller nodes run Raft consensus over an internal replicated log topic named \textbf{\texttt{\_\_cluster\_metadata}}, enabling millions of partitions and sub-second controller failovers.

\subsection{Producer Batching and Kafka 3.0+ Idempotent Defaults}

Sending one network packet per $200\text{-byte}$ JSON event wastes TCP overhead. Instead, Kafka producers accumulate records in RAM (\texttt{buffer.memory}, default $32\text{ MB}$) into per-partition batches. A sender thread dispatches a batch as soon as \emph{either}:
\begin{enumerate}[leftmargin=18pt]
  \item The batch reaches \texttt{batch.size} (default $16{,}384$ bytes), \textbf{or}
  \item \texttt{linger.ms} expires (default $0\text{ ms}$; setting $10\text{--}20\text{ ms}$ lets batches fill cleanly under load).
\end{enumerate}
If records sit in the accumulator longer than \texttt{delivery.timeout.ms} (default $120{,}000\text{ ms}$) due to downstream slowness or uncompressed payloads saturating bandwidth, the producer raises \texttt{TimeoutException}.

Since Kafka~3.0, producers enable \texttt{acks=all}, \texttt{enable.idempotence=true}, and \texttt{max.in.flight.requests.per.connection=5} by default. Why is \texttt{5} safe? Without idempotence, allowing 5 unacknowledged requests on the wire risks reordering if batch~1 fails and retries behind batch~2. With idempotence enabled, the producer stamps each batch with a monotonically increasing sequence number; the broker rejects out-of-order or duplicate sequences, giving both high pipelined throughput and strict per-partition ordering!

\subsection{Consumer Groups and Rebalance Strategies}

Within a single \textbf{Consumer Group}, Kafka guarantees that each partition is assigned to \textbf{at most one} consumer instance at a time (if you have 3 partitions and 4 workers, the 4th worker stays idle). When a worker joins, crashes, or misses its heartbeat, the group coordinator runs a partition reassignment (\textbf{rebalance}):
\begin{itemize}[leftmargin=18pt]
  \item \textbf{\texttt{RangeAssignor}} (eager): Assigns contiguous partition slices per topic (e.g., Topic~A partitions 0--1 to C1, 2 to C2). On rebalance, every consumer revokes all partitions and stops reading until reassignment completes (a stop-the-world stall).
  \item \textbf{\texttt{RoundRobinAssignor}} (eager): Stripes partitions across consumers evenly across all subscribed topics, preventing skew when topics have few partitions, but still halts all consumers during rebalance.
  \item \textbf{\texttt{CooperativeStickyAssignor}} (incremental): Keeps existing assignments intact and revokes \emph{only} the specific partitions that need to move to a new or surviving consumer. Unaffected partitions keep streaming with zero downtime.
\end{itemize}

\begin{lstlisting}[language=Python, caption={Kafka Idempotent Producer and Cooperative Consumer in Python (\texttt{confluent\_kafka})}]
import json
from confluent_kafka import Producer, Consumer

# 1. Producer with batching, LZ4 compression, and Kafka 3.0+ safety defaults
producer = Producer({
    "bootstrap.servers": "kafka-1:9092,kafka-2:9092",
    "acks": "all",
    "enable.idempotence": True,
    "max.in.flight.requests.per.connection": 5,
    "compression.type": "lz4",
    "batch.size": 65536,   # 64 KB batches
    "linger.ms": 15,       # Wait up to 15ms to fill batch
})

order = {"order_id": "ord-501", "account_id": "acct-9", "amount": 120.0}
producer.produce(
    topic="orders",
    key=order["account_id"].encode("utf-8"),
    value=json.dumps(order).encode("utf-8"),
)
producer.flush(timeout=5.0)

# 2. Consumer with CooperativeStickyAssignor and manual offset commit
consumer = Consumer({
    "bootstrap.servers": "kafka-1:9092,kafka-2:9092",
    "group.id": "billing-workers",
    "auto.offset.reset": "earliest",
    "enable.auto.commit": False,
    "partition.assignment.strategy": "cooperative-sticky",
})
consumer.subscribe(["orders"])

msg = consumer.poll(timeout=1.0)
if msg is not None and not msg.error():
    payload = json.loads(msg.value().decode("utf-8"))
    # Process payload idempotently in DB first, then commit offset
    consumer.commit(message=msg, asynchronous=False)
\end{lstlisting}

\textbf{Step-by-step walkthrough:}
\begin{enumerate}[leftmargin=18pt]
  \item Lines~5--13 configure the producer to require ISR quorum confirmation (\texttt{acks="all"}), prevent retry duplicates/reordering (\texttt{enable.idempotence=True} with \texttt{max.in.flight=5}), compress batches with LZ4, and buffer up to $64\text{ KB}$ or $15\text{ ms}$.
  \item Lines~15--21 key the message by \texttt{"acct-9"}. Kafka hashes the key (\texttt{murmur2(key) \% num\_partitions}) so all events for \texttt{acct-9} always land on the same partition in strict chronological order; \texttt{flush()} blocks until the broker ISR acknowledges.
  \item Lines~24--32 create a consumer in group \texttt{"billing-workers"} disabling auto-commit (so offsets never advance before DB writes succeed) and enabling \texttt{"cooperative-sticky"} incremental rebalancing.
  \item Lines~34--38 poll one record, decode the JSON order, and call synchronous \texttt{commit(message=msg)} only \emph{after} processing completes, enforcing at-least-once semantics.
\end{enumerate}

\section{RabbitMQ Internals}

Written in Erlang/OTP and implementing the \textbf{AMQP 0-9-1} protocol, RabbitMQ excels at smart routing, task prioritization, and per-message acknowledgments rather than multi-day log retention.

\subsection{AMQP 0-9-1 Exchanges and Bindings}

In RabbitMQ, producers never publish directly to queues. They publish to an \textbf{Exchange}, which evaluates rules called \textbf{Bindings} to route copies into one or more \textbf{Queues}:
\begin{enumerate}[leftmargin=18pt]
  \item \textbf{Direct}: Routes when the message's routing key equals the binding key exactly.
  \item \textbf{Fanout}: Ignores routing keys; broadcasts copies to every bound queue.
  \item \textbf{Topic}: Matches dot-separated words using \texttt{*} (exactly one word) and \texttt{\#} (zero or more words).
  \item \textbf{Headers}: Routes using arbitrary key-value attributes in message headers.
\end{enumerate}

\begin{figure}[htbp]
\centering
\begin{tikzpicture}[
  scale=0.85,
  p/.style={rectangle, draw=headingblue, fill=codebg, thick, rounded corners=2pt, minimum width=2.0cm, minimum height=1.0cm, align=center, font=\sffamily\footnotesize},
  ex/.style={circle, draw=exampleBlue, fill=exampleBg, thick, minimum size=1.8cm, align=center, font=\sffamily\footnotesize},
  q/.style={rectangle, draw=intuitionTeal, fill=intuitionBg, thick, rounded corners=2pt, minimum width=2.4cm, minimum height=0.9cm, align=center, font=\sffamily\footnotesize}
]
  \node[p] (prod) at (0, 0) {Producer};
  \node[ex] (exch) at (4, 0) {\textbf{Topic}\\Exchange};
  \node[q] (q1) at (9.2, 1.5) {Queue A\\(\texttt{order.*})};
  \node[q] (q2) at (9.2, -1.5) {Queue B\\(\texttt{*.eu})};

  \draw[-{Stealth[scale=1.0]}, thick] (prod) -- node[above, font=\sffamily\tiny] {publish(\texttt{order.eu})} (exch);
  \draw[-{Stealth[scale=1.0]}, thick, color=intuitionTeal] (exch) -- node[above, sloped, font=\sffamily\tiny] {Matches \texttt{order.*}} (q1);
  \draw[-{Stealth[scale=1.0]}, thick, color=intuitionTeal] (exch) -- node[below, sloped, font=\sffamily\tiny] {Matches \texttt{*.eu}} (q2);
\end{tikzpicture}
\caption{RabbitMQ AMQP 0-9-1 Topic Exchange Routing Topology}
\label{fig:v4-ch05-amqp-routing}
\end{figure}

\begin{lstlisting}[language=Python, caption={RabbitMQ Topic Exchange, Prefetch, and Manual Ack in Python (\texttt{pika})}]
import json
import pika

conn = pika.BlockingConnection(pika.ConnectionParameters("rabbitmq-1"))
ch = conn.channel()

# Declare a durable topic exchange and queue, then bind with wildcard pattern
ch.exchange_declare(exchange="orders_ex", exchange_type="topic", durable=True)
ch.queue_declare(queue="eu_orders", durable=True, arguments={"x-queue-type": "quorum"})
ch.queue_bind(queue="eu_orders", exchange="orders_ex", routing_key="order.*.eu")

# Publish a persistent message matching "order.created.eu"
ch.basic_publish(
    exchange="orders_ex",
    routing_key="order.created.eu",
    body=json.dumps({"order_id": "ord-501"}).encode("utf-8"),
    properties=pika.BasicProperties(delivery_mode=2),  # Persist to disk
)

# Consume safely with bounded prefetch backpressure
ch.basic_qos(prefetch_count=50)
method, props, body = ch.basic_get(queue="eu_orders", auto_ack=False)
if method is not None:
    data = json.loads(body.decode("utf-8"))
    ch.basic_ack(delivery_tag=method.delivery_tag)
conn.close()
\end{lstlisting}

\textbf{Step-by-step walkthrough:}
\begin{enumerate}[leftmargin=18pt]
  \item Lines~8--10 declare a durable topic exchange \texttt{orders\_ex}, a Raft-backed \texttt{quorum} queue \texttt{eu\_orders}, and a binding \texttt{order.*.eu} (note \texttt{*} matches one word, \texttt{created}).
  \item Lines~13--18 publish \texttt{\{"order\_id": "ord-501"\}} with routing key \texttt{order.created.eu} and \texttt{delivery\_mode=2} (persistent on disk). The exchange matches \texttt{order.*.eu} and routes the message into \texttt{eu\_orders}.
  \item Lines~21--25 bound unacked deliveries to 50 (\texttt{basic\_qos}), fetch the message with \texttt{auto\_ack=False}, parse the JSON, and send \texttt{basic\_ack(delivery\_tag)} so the queue deletes the message only after completion.
\end{enumerate}

\subsection{Erlang BEAM Reductions, Mnesia, and Quorum Queues}

RabbitMQ runs on the Erlang **BEAM** virtual machine, where every TCP connection, AMQP channel, and queue is an isolated lightweight **actor process** ($\sim 2.5\text{ KB}$ memory, sharing no state and communicating via async messages).
\begin{itemize}[leftmargin=18pt]
  \item \textbf{BEAM Reductions}: Rather than letting one busy queue monopolize a CPU core, the BEAM scheduler tracks work using \textbf{reductions} (an abstract counter roughly equal to one function call or built-in operation). After **4,000 reductions**, BEAM preempts the actor and schedules the next waiting actor---ensuring smooth tail latency across hundreds of thousands of queues.
  \item \textbf{Mnesia}: Erlang's built-in distributed database stores cluster metadata (users, vhosts, exchanges, queue definitions, and binding rules) in RAM and on disk across broker nodes.
  \item \textbf{Classic Mirrored Queues vs. Quorum Queues}: Legacy Classic Mirrored Queues used primary/mirror replication that silently dropped un-synced messages on crash and froze traffic during resynchronization. Modern **Quorum Queues** (RabbitMQ 3.8+, set via \texttt{x-queue-type: quorum}) implement **Raft consensus** per queue across an odd number of brokers ($3$ or $5$), surviving crashes with zero message loss and non-blocking catch-up.
\end{itemize}

\begin{realworldbox}[Diagnosing a Kafka Consumer Rebalance Storm]
\textbf{Symptom:} A \texttt{billing-workers} deployment repeatedly logs \texttt{CommitFailedException} and consumer lag climbs to 500,000 messages while CPU sits below 20\%.\\
\textbf{Sample Broker / Client Logs:}
\begin{verbatim}
[Consumer group=billing-workers] Member consumer-3-a81f exceeded max.poll.interval.ms
(300000 ms) and was kicked out of the group. Triggering rebalance (generation 47 -> 48).
CommitFailedException: Commit cannot be completed since the group has already rebalanced.
\end{verbatim}
\textbf{Root Cause \& Fix:} A slow downstream API call stalled a worker's loop for $>5\text{ minutes}$ between \texttt{poll()} calls, causing the coordinator to evict the worker and trigger an eager rebalance storm across the whole group. Lower \texttt{max.poll.records} (e.g., from 500 to 50) so each batch finishes well inside \texttt{max.poll.interval.ms}, switch to \texttt{CooperativeStickyAssignor}, and offload slow calls or set strict HTTP timeouts.
\end{realworldbox}

\begin{tryityourselfbox}[Trace Partition Assignment and AMQP Wildcards]
\begin{enumerate}[leftmargin=18pt]
  \item Suppose Topic~T1 has 4 partitions ($P_0..P_3$) and Topic~T2 has 4 partitions ($P_0..P_3$). A consumer group has 2 workers ($C_1, C_2$) subscribed to both topics. List the exact partition assignment under \texttt{RangeAssignor} vs. \texttt{RoundRobinAssignor}. What happens if a 3rd worker $C_3$ joins under \texttt{RangeAssignor}?
  \item Which of the following AMQP routing keys match a binding of \texttt{payments.\#.failed}: (a) \texttt{payments.failed}, (b) \texttt{payments.eu.card.failed}, (c) \texttt{payments.card.succeeded}?
\end{enumerate}
\end{tryityourselfbox}

\section{Summary \protect\& Key Takeaways}
\begin{itemize}[leftmargin=18pt]
  \item Kafka pairs sequential disk appends, sparse \texttt{.index} binary search, the Linux OS Page Cache, and \texttt{sendfile(2)} zero-copy (2 DMA copies, 0 CPU copies) to deliver multi-GB/s throughput without JVM heap bloat.
  \item Consumers only read up to the High Watermark ($\min(\text{LEO}_{\text{ISR}})$); KRaft replaces ZooKeeper with an internal Raft topic named \texttt{\_\_cluster\_metadata}.
  \item Kafka 3.0+ defaults (\texttt{acks=all}, \texttt{enable.idempotence=true}, \texttt{max.in.flight=5}) plus \texttt{CooperativeStickyAssignor} maximize throughput without sacrificing order or triggering stop-the-world rebalances.
  \item RabbitMQ routes messages through AMQP Exchanges (Direct, Fanout, Topic, Headers), stores topology in Mnesia, schedules actors preemptively via BEAM reductions (4,000 per slice), and achieves Raft-backed HA via Quorum Queues.
\end{itemize}

\section{Exercises}

\subsection*{Exercise~\thechapter.1 (Mechanics): Partition Count vs. Consumer Scaling}
Given a Kafka topic with 3 partitions and a consumer group with 4 active workers:
\begin{enumerate}[leftmargin=18pt]
  \item How many workers actively consume records under \texttt{RangeAssignor}?
  \item What happens if one active worker crashes?
  \item Why does \texttt{CooperativeStickyAssignor} prevent full-group processing stalls during that failover?
\end{enumerate}

\subsection*{Exercise~\thechapter.2 (Debugging): Kafka Producer \texttt{TimeoutException}}
A telemetry producer throws:
\begin{lstlisting}[language=Java]
org.apache.kafka.common.errors.TimeoutException: Expiring 48 record(s) for sensors-2: 120000 ms has passed since batch creation
\end{lstlisting}
Network links are healthy and broker CPU is 15\%. Explain how \texttt{buffer.memory}, \texttt{batch.size}, \texttt{linger.ms}, \texttt{compression.type}, and \texttt{delivery.timeout.ms} interact to cause and fix this error.

\subsection*{Exercise~\thechapter.3 (Architecture): Selecting Kafka vs. RabbitMQ}
Justify choosing Kafka or RabbitMQ for:
\begin{enumerate}[leftmargin=18pt]
  \item A core ledger event backbone needing 7-year replayable audit history, strict per-account ordering, and fraud-detection stream joins.
  \item A multi-step PDF/video rendering dispatcher needing VIP priority queues, wildcard routing to GPU/CPU pools, and immediate deletion upon task completion.
\end{enumerate}

\subsection*{Exercise~\thechapter.4 (Try-It-Yourself Lab): Assignor Trace & AMQP Topic Matching}
Solve both parts of the \emph{Try It Yourself} box above.

\section{Detailed Solutions}

\subsection*{Solution to Exercise~\thechapter.1}
\begin{enumerate}[leftmargin=18pt]
  \item At most one worker per partition in a group means \textbf{3 workers} consume partitions ($P_0, P_1, P_2$) while the 4th worker remains idle standby.
  \item On crash, the coordinator detects missed heartbeats and reassigns the orphaned partition to the idle 4th worker.
  \item Eager assignors revoke \emph{all} partitions from all workers before reassigning. \texttt{CooperativeStickyAssignor} keeps unaffected partitions running continuously and migrates only the single orphaned partition.
\end{enumerate}

\subsection*{Solution to Exercise~\thechapter.2}
Uncompressed JSON telemetry fills \texttt{buffer.memory} faster than the sender thread can drain tiny $16\text{ KB}$ batches, leaving queued batches waiting past \texttt{delivery.timeout.ms} ($120\text{ s}$). Fix by enabling \texttt{compression.type=lz4} (shrinking payloads 60--80\%), raising \texttt{batch.size=65536}, setting \texttt{linger.ms=20}, and keeping \texttt{max.in.flight.requests.per.connection=5} with \texttt{enable.idempotence=true}.

\subsection*{Solution to Exercise~\thechapter.3}
\begin{itemize}[leftmargin=18pt]
  \item \textbf{Core Ledger $\to$ Apache Kafka}: Tiered append-only storage retains 7 years of immutable events; partitioning by \texttt{account\_id} guarantees strict per-account order; Kafka Streams/Flink processes sliding fraud windows natively.
  \item \textbf{Rendering Dispatcher $\to$ RabbitMQ}: AMQP Topic exchanges route \texttt{render.gpu.*} vs. \texttt{render.cpu.*}; native priority queues (\texttt{x-max-priority}) let VIP jobs jump queues; destructive ACKs purge finished jobs immediately.
\end{itemize}

\subsection*{Solution to Exercise~\thechapter.4}
\begin{enumerate}[leftmargin=18pt]
  \item Under \texttt{RangeAssignor} (assigned per topic): $C_1$ gets \texttt{T1-0, T1-1, T2-0, T2-1}; $C_2$ gets \texttt{T1-2, T1-3, T2-2, T2-3}. Under \texttt{RoundRobinAssignor}: identical 4/4 split here, though on 3 consumers ($C_1, C_2, C_3$) \texttt{RangeAssignor} skews heavily: $C_1$ gets 2 partitions per topic (\texttt{T1-0, T1-1, T2-0, T2-1} = 4 total) while $C_2$ and $C_3$ get only 2 total each! \texttt{RoundRobinAssignor} gives a balanced 3/3/2 split.
  \item Since \texttt{\#} matches zero or more words: (a) \texttt{payments.failed} matches (0 words between), (b) \texttt{payments.eu.card.failed} matches (2 words between), (c) \texttt{payments.card.succeeded} does \emph{not} match.
\end{enumerate}

\begin{labbridgebox}[Event-Driven Outbox & Full Stack Broker Integration]
Explore Kafka-style partitioned log delivery in \texttt{projects/p11-event-driven-microservices}, and inspect live RabbitMQ exchanges, queues, and worker consumption in \texttt{projects/p12-end-to-end-platform} (ports 5672/15672).
\end{labbridgebox}

\section{Further Reading \protect\& Citations}
\begin{itemize}[leftmargin=18pt]
  \item Kreps, J., Narkhede, N., \& Rao, J. (2011). \emph{Kafka: A distributed messaging system for log processing}. NetDB Workshop, 1--7.
  \item Videla, A., \& Williams, J. (2012). \emph{RabbitMQ in Action: Distributed messaging for everyone}. Manning Publications.
  \item Armstrong, J. (2010). \emph{Erlang}. Communications of the ACM, 53(9), 68--75.
\end{itemize}
